\documentclass[11pt,a4paper]{article}

\usepackage[utf8]{inputenc}
\usepackage[T1]{fontenc}
\usepackage[french]{babel}
\usepackage{geometry}
\usepackage{hyperref}
\usepackage{longtable}
\usepackage{array}
\usepackage{enumitem}
\usepackage{xcolor}
\usepackage{listings}
\usepackage{fancyhdr}

\geometry{margin=2.3cm}
\hypersetup{
  colorlinks=true,
  linkcolor=blue,
  urlcolor=blue,
  pdftitle={ARC-MEM Bridge v0.1 - Architecture et tutoriel},
  pdfauthor={ARC-MEM Bridge}
}

\pagestyle{fancy}
\fancyhf{}
\lhead{ARC-MEM Bridge v0.1}
\rhead{ZORA\_ARC V9.2 ARC-MEM-H}
\cfoot{\thepage}

\lstdefinestyle{console}{
  basicstyle=\ttfamily\small,
  breaklines=true,
  frame=single,
  backgroundcolor=\color{gray!6},
  columns=fullflexible
}
\lstset{style=console}

\title{\textbf{ARC-MEM Bridge v0.1}\\Architecture du système et tutoriel de fonctionnement}
\author{Architecture cible : ZORA\_ARC V9.2 ARC-MEM-H}
\date{13 juin 2026}

\begin{document}
\maketitle

\begin{abstract}
ARC-MEM Bridge est une couche locale d'intégration entre Hermes Desktop, Hindsight,
MemGraphRAG, H-MEM, PostgreSQL/PostGIS/pgvector et un moteur Python FastAPI. Ce
document décrit l'architecture, la séparation des responsabilités, les flux de
fonctionnement, les schémas PostgreSQL, les commandes d'exploitation et un tutoriel
pratique pour diagnostiquer, démarrer, ingérer, récupérer et auditer la mémoire.
\end{abstract}

\tableofcontents
\newpage

\section{Objectif du système}

ARC-MEM Bridge sert de pont logique local entre les composants de mémoire et
d'orchestration. Il ne remplace ni Hermes, ni Hindsight, ni MemGraphRAG. Il fournit
une API locale, une CLI et des services Python pour unifier l'accès aux mémoires,
contrôler la provenance et appliquer la gouvernance ZORA ARC.

\subsection{Objectifs principaux}

\begin{itemize}[leftmargin=*]
  \item connecter proprement PostgreSQL depuis WSL ;
  \item garder Hindsight comme mémoire agentique persistante ;
  \item garder MemGraphRAG comme mémoire documentaire probatoire ;
  \item utiliser H-MEM comme routeur hiérarchique, sans base séparée ;
  \item exposer une API locale pour Hermes ;
  \item appliquer des contrôles ZORA ARC : TASKCARD, preuve, audit, gates ;
  \item éviter toute fuite de secrets dans les logs, tests et documentation.
\end{itemize}

\section{Vue d'ensemble}

\begin{verbatim}
                 +----------------------+
                 |    Hermes Desktop    |
                 | Interface agentique  |
                 +----------+-----------+
                            |
                            | API locale / CLI
                            v
                 +----------------------+
                 |   ARC-MEM Bridge     |
                 | FastAPI + services   |
                 +----+------+-----+----+
                      |      |     |
      +---------------+      |     +----------------+
      |                      |                      |
      v                      v                      v
+-------------+     +----------------+      +----------------+
| Hindsight   |     | MemGraphRAG    |      | H-MEM Router   |
| agentique   |     | documentaire   |      | hierarchique   |
| PostgreSQL  |     | probatoire     |      | chemins memoire|
+------+------+     +-------+--------+      +-------+--------+
       |                    |                       |
       +--------------------+-----------------------+
                            |
                            v
                +------------------------+
                | PostgreSQL Windows 18  |
                | PostGIS + pgvector     |
                +------------------------+
\end{verbatim}

\section{Etat cible retenu}

L'installation actuelle retient le mode \texttt{postgresql\_direct} pour Hindsight.
Ce choix abandonne explicitement PostgreSQL embedded, qui est fragile sur Windows
avec des locales regionales non ASCII.

\begin{longtable}{|p{3.2cm}|p{11cm}|}
\hline
\textbf{Composant} & \textbf{Etat cible} \\
\hline
Hermes Desktop & Interface agent. Ne doit pas etre modifie brutalement. Les
fichiers de configuration doivent etre sauvegardes avant toute modification. \\
\hline
Hindsight & Memoire agentique persistante en mode \texttt{postgresql\_direct}.
Stocke preferences, decisions, contexte de travail et resumes intersessions dans
\texttt{hindsight\_hermes}. \\
\hline
MemGraphRAG & Memoire documentaire probatoire dans les schemas ARC-MEM. Stocke
documents, passages, entites, faits, relations, conflits et preuves. \\
\hline
H-MEM & Couche de routage hierarchique. Les niveaux sont \texttt{domain},
\texttt{theme}, \texttt{subtheme}, \texttt{entity}, \texttt{fact},
\texttt{passage}. \\
\hline
PostgreSQL & Serveur central Windows 18 accessible depuis WSL via
\texttt{172.21.96.1:5432}. Extensions actives : \texttt{vector},
\texttt{postgis}, \texttt{pg\_trgm}, \texttt{pgcrypto}. \\
\hline
ZORA ARC & Gouvernance methodologique : TASKCARD, score routeur, registre de
preuves, audit, niveaux de criticite. \\
\hline
\end{longtable}

\section{Separation stricte des responsabilites}

\subsection{Hindsight}

Hindsight reste reserve a la memoire agentique persistante :

\begin{itemize}[leftmargin=*]
  \item preferences utilisateur ;
  \item decisions de projet ;
  \item historique de travail ;
  \item resumes intersessions ;
  \item contexte Hermes ;
  \item faits agentiques courts et recherchables.
\end{itemize}

Il ne doit pas recevoir les PDF complets ni les gros chunks documentaires. Dans
l'etat actuel, Hindsight est connecte directement a PostgreSQL dans la base
\texttt{hindsight\_hermes}. Le daemon Hindsight et PostgreSQL embedded ne sont
pas necessaires pour cette voie.

\subsection{MemGraphRAG}

MemGraphRAG est reserve a la memoire documentaire probatoire :

\begin{itemize}[leftmargin=*]
  \item documents sources ;
  \item chunks ;
  \item passages ;
  \item entites ;
  \item faits documentaires ;
  \item conflits ;
  \item preuves et provenance.
\end{itemize}

La chaine de tracabilite visee est :

\begin{center}
\texttt{source -> chunk -> passage -> fait -> citation -> reponse}
\end{center}

\subsection{H-MEM}

H-MEM ne cree pas une base de memoire parallele. Il route les requetes vers les
bons niveaux d'abstraction et les bons objets. Les noeuds H-MEM peuvent pointer
vers :

\begin{itemize}[leftmargin=*]
  \item \texttt{zora.documents} ;
  \item \texttt{zora.chunks} ;
  \item \texttt{memgraph.mem\_entities} ;
  \item \texttt{memgraph.mem\_facts} ;
  \item \texttt{memgraph.mem\_passages} ;
  \item des souvenirs Hindsight par reference externe.
\end{itemize}

\section{Organisation PostgreSQL}

\subsection{Bases}

\begin{longtable}{|p{4cm}|p{10cm}|}
\hline
\textbf{Base} & \textbf{Role} \\
\hline
\texttt{rag\_arc} & Base ARC-MEM Bridge. Contient les schemas
\texttt{zora}, \texttt{memgraph}, \texttt{hmem}. \\
\hline
\texttt{hindsight\_hermes} & Base Hindsight directe. Contient la memoire
agentique deja migree : faits, sessions, entites, documents, memories. \\
\hline
\end{longtable}

\subsection{Schemas ARC-MEM dans \texttt{rag\_arc}}

\begin{longtable}{|p{3cm}|p{11cm}|}
\hline
\textbf{Schema} & \textbf{Tables principales} \\
\hline
\texttt{zora} & \texttt{taskcards}, \texttt{arc\_runs},
\texttt{arc\_audit\_log}, \texttt{evidence\_registry},
\texttt{documents}, \texttt{chunks}, \texttt{citations}. \\
\hline
\texttt{memgraph} & \texttt{mem\_schemas}, \texttt{mem\_entities},
\texttt{mem\_facts}, \texttt{mem\_passages}, \texttt{mem\_bridges},
\texttt{mem\_conflicts}, \texttt{v\_graphe\_ppr}. \\
\hline
\texttt{hmem} & \texttt{memory\_layers}, \texttt{memory\_nodes},
\texttt{memory\_edges}, \texttt{positional\_indexes},
\texttt{routing\_paths}. \\
\hline
\end{longtable}

\section{Application Python}

\subsection{Structure}

\begin{lstlisting}
arc-mem-bridge/
  app/
    main.py
    config.py
    db.py
    api/
      health.py
      memory.py
      ingest.py
      retrieve.py
      audit.py
    services/
      hindsight_client.py
      memgraph_service.py
      hmem_router.py
      zora_arc_service.py
      retrieval_service.py
      ppr_service.py
      hermes_bridge.py
    migrations/
      001_init_extensions.sql
      002_init_zora_schema.sql
      003_init_memgraph_schema.sql
      004_init_hmem_schema.sql
      005_views_and_indexes.sql
    cli/
      arc_mem.py
  scripts/
    check_postgres.py
    apply_migrations.py
    diagnose_hermes.py
    diagnose_hindsight.py
    smoke_test.py
\end{lstlisting}

\subsection{Routes API}

\begin{longtable}{|p{4cm}|p{10cm}|}
\hline
\textbf{Route} & \textbf{Role} \\
\hline
\texttt{GET /health} & Retourne le statut de l'application, PostgreSQL,
extensions et Hindsight. \\
\hline
\texttt{POST /memory/retain} & Retient un souvenir agentique dans Hindsight.
En mode direct, insere dans \texttt{hindsight\_hermes.facts}. \\
\hline
\texttt{POST /memory/recall} & Recherche des souvenirs agentiques Hindsight. \\
\hline
\texttt{POST /ingest/document} & Enregistre un document dans \texttt{zora},
le decoupe en chunks et prepare l'extraction probatoire. \\
\hline
\texttt{POST /retrieve} & Construit un paquet de contexte structure :
TASKCARD, contexte Hindsight, chemins H-MEM, faits, passages, citations. \\
\hline
\texttt{POST /audit} & Verifie qu'une reponse proposee contient sources,
citations, limites, risques et coherence avec la TASKCARD. \\
\hline
\end{longtable}

\section{Tutoriel de fonctionnement}

\subsection{1. Activer l'environnement Python}

\begin{lstlisting}
cd /mnt/c/Users/Stephane_Arnoux/arc-mem-bridge
. .venv/bin/activate
\end{lstlisting}

\subsection{2. Verifier PostgreSQL}

\begin{lstlisting}
python scripts/check_postgres.py
\end{lstlisting}

Resultat attendu :

\begin{lstlisting}
"ok": true
"database": "rag_arc"
"vector": true
"postgis": true
"pg_trgm": true
"pgcrypto": true
\end{lstlisting}

\subsection{3. Verifier Hindsight en PostgreSQL direct}

\begin{lstlisting}
python scripts/diagnose_hindsight.py
\end{lstlisting}

Resultat attendu :

\begin{lstlisting}
"ok": true
"mode": "postgresql_direct"
"facts": 661
"sessions": 212
"entities": 14
\end{lstlisting}

\subsection{4. Appliquer ou verifier les migrations}

Les migrations sont idempotentes. Pour les auditer sans execution :

\begin{lstlisting}
python scripts/apply_migrations.py --dry-run
\end{lstlisting}

Pour les appliquer :

\begin{lstlisting}
python scripts/apply_migrations.py
\end{lstlisting}

\subsection{5. Demarrer l'API locale}

\begin{lstlisting}
uvicorn app.main:app --host 127.0.0.1 --port 8791
\end{lstlisting}

Puis tester :

\begin{lstlisting}
curl http://127.0.0.1:8791/health
\end{lstlisting}

\subsection{6. Utiliser la CLI}

\begin{lstlisting}
python -m app.cli.arc_mem health
python -m app.cli.arc_mem check-postgres
python -m app.cli.arc_mem diagnose-hermes
python -m app.cli.arc_mem diagnose-hindsight
python -m app.cli.arc_mem smoke-test
\end{lstlisting}

\subsection{7. Retenir un souvenir agentique}

\begin{lstlisting}
curl -X POST http://127.0.0.1:8791/memory/retain \
  -H "Content-Type: application/json" \
  -d '{
    "text": "ARC-MEM Bridge utilise Hindsight en postgresql_direct.",
    "metadata": {"source": "tutorial", "type": "configuration"}
  }'
\end{lstlisting}

En mode \texttt{postgresql\_direct}, cette route insere une ligne dans
\texttt{hindsight\_hermes.facts}. Elle ne duplique pas ce souvenir dans
MemGraphRAG.

\subsection{8. Rappeler un souvenir agentique}

\begin{lstlisting}
curl -X POST http://127.0.0.1:8791/memory/recall \
  -H "Content-Type: application/json" \
  -d '{"query": "postgresql_direct", "limit": 5}'
\end{lstlisting}

\subsection{9. Ingerer un document probatoire}

\begin{lstlisting}
curl -X POST http://127.0.0.1:8791/ingest/document \
  -H "Content-Type: application/json" \
  -d '{
    "title": "Note architecture ARC-MEM",
    "content": "MemGraphRAG conserve les passages et les faits documentaires.",
    "source_uri": "local://tutorial",
    "metadata": {"domain": "architecture"}
  }'
\end{lstlisting}

Cette operation ecrit dans \texttt{zora.documents} et \texttt{zora.chunks}. Les
embeddings et l'extraction LLM restent des etapes configurees separement.

\subsection{10. Recuperer un paquet de contexte}

\begin{lstlisting}
curl -X POST http://127.0.0.1:8791/retrieve \
  -H "Content-Type: application/json" \
  -d '{
    "query": "Comment ARC-MEM separe Hindsight et MemGraphRAG ?",
    "expected_evidence_level": "standard",
    "limit": 8
  }'
\end{lstlisting}

La reponse n'est pas une reponse finale longue. Elle contient un
\texttt{ContextBundle} :

\begin{itemize}[leftmargin=*]
  \item \texttt{memories} : contexte agentique Hindsight ;
  \item \texttt{route\_paths} : chemins H-MEM ;
  \item \texttt{facts} : faits MemGraphRAG ;
  \item \texttt{passages} : chunks documentaires ;
  \item \texttt{citations} : provenance ;
  \item \texttt{warnings} : limites et indisponibilites.
\end{itemize}

\subsection{11. Auditer une reponse}

\begin{lstlisting}
curl -X POST http://127.0.0.1:8791/audit \
  -H "Content-Type: application/json" \
  -d '{
    "taskcard": {"title": "Audit exemple"},
    "proposed_answer": "Reponse courte",
    "sources": [],
    "citations": [],
    "limits": [],
    "risks": []
  }'
\end{lstlisting}

Si sources, citations, limites ou risques manquent, l'audit retourne des
\texttt{findings}.

\section{Flux de retrieval}

\begin{enumerate}[leftmargin=*]
  \item L'utilisateur ou Hermes appelle \texttt{POST /retrieve}.
  \item ZORA ARC cree une TASKCARD minimale.
  \item Le routeur ZORA ARC calcule les signaux \texttt{Ct}, \texttt{Cs},
  \texttt{Rk}, \texttt{Ex}, \texttt{Op}, \texttt{Tp}.
  \item Hindsight fournit le contexte agentique via PostgreSQL direct.
  \item H-MEM choisit un chemin de routage hierarchique.
  \item MemGraphRAG recherche les faits/passages probatoires.
  \item Le service de retrieval assemble un \texttt{ContextBundle}.
\end{enumerate}

\section{Securite et secrets}

\begin{itemize}[leftmargin=*]
  \item Ne jamais exposer \texttt{.env}.
  \item Ne jamais afficher de cle complete dans les logs.
  \item Les motifs \texttt{nvapi-*}, \texttt{sk-*}, tokens, passwords et secrets
  sont masques par la couche de logging.
  \item Les configurations Hermes/Hindsight sont sauvegardees avant modification.
  \item PostgreSQL embedded est abandonne pour eviter les erreurs de locale
  Windows et les locks de fichiers.
\end{itemize}

\section{Commandes de validation}

\begin{lstlisting}
python scripts/check_postgres.py
python scripts/diagnose_hindsight.py
python scripts/apply_migrations.py --dry-run
python -m pytest
python -m app.cli.arc_mem health
uvicorn app.main:app --host 127.0.0.1 --port 8791
\end{lstlisting}

\section{Etat valide au 13 juin 2026}

\begin{longtable}{|p{5cm}|p{9cm}|}
\hline
\textbf{Controle} & \textbf{Resultat} \\
\hline
PostgreSQL \texttt{rag\_arc} & OK \\
\hline
Extensions & \texttt{vector}, \texttt{postgis}, \texttt{pg\_trgm},
\texttt{pgcrypto} OK \\
\hline
Hindsight & OK en \texttt{postgresql\_direct} \\
\hline
Memoire Hindsight & 661 faits, 212 sessions, 14 entites \\
\hline
Tests Python & 8 tests passent \\
\hline
API locale & \texttt{/health} repond OK \\
\hline
\end{longtable}

\end{document}
